\chapter{Temporomandibular joint dysfunction}
\index{Temporomandibular joint dysfunction}

\medskip
\noindent\textit{Educational reference; always seek professional advice for individual care.}

\section*{Definition and Overview}
Temporomandibular joint dysfunction, often shortened to TMD or TMJ dysfunction, is a group of conditions that cause pain and impaired function in the jaw joint and the muscles that move it. The temporomandibular joint sits just in front of each ear, where the lower jaw (mandible) meets the skull, and it is one of the most used joints in the body, opening and closing thousands of times a day with every bite, word, and yawn. TMD is not one disease but a family of problems, including muscle spasm and tenderness, inflammation of the joint, displacement of the small disc that cushions it, and arthritis of the joint itself. People with TMD typically have some combination of jaw pain, clicking or grinding noises, difficulty opening the mouth, and pain in the face, ears, or head. The condition is common, affects women more often than men, and is frequently driven or worsened by clenching and grinding of the teeth, often during sleep and often during stressful periods. The overwhelming majority of cases are managed successfully with simple, conservative measures: education, rest, a soft diet, heat, gentle exercises, stress management, and short-term pain relief. This chapter explains what TMD is, how it is recognised, which conditions can mimic it, and what treatment actually works, and it emphasises the reassuring reality that surgery is rarely needed.

\section*{Epidemiology}
TMD is one of the commonest causes of facial pain, affecting a substantial minority of the population at some point in life. Surveys suggest that around one person in five has some symptoms of TMD, although only a smaller proportion ever seeks treatment, and women are affected roughly twice as often as men, with symptoms most likely to begin between the ages of twenty and forty. The condition is less common in children, and in older adults symptoms are more likely to reflect arthritis of the joint. Mild clicking of the jaw is very common and often harmless; many people have occasional joint noises without pain or limitation and never need treatment. The frequency of TMD appears to have been stable over recent decades, and, despite older ideas, there is little evidence that routine orthodontic treatment or wisdom tooth removal causes it. Stress, poor sleep, and anxiety are associated with flare-ups, and people who clench or grind their teeth, particularly at night, are over-represented among those who seek help.

\section*{Aetiology and Pathophysiology}
The jaw joint is a complex structure: a shallow socket in the skull, the rounded condyle of the lower jaw, and a small cushioning disc between them, held together by ligaments and powered by several muscles of chewing. Pain and dysfunction arise when any part of this system is overloaded or inflamed.

\begin{itemize}
    \item \textbf{Muscle overactivity:} clenching and grinding (bruxism), often during sleep or stress, overload the jaw muscles and cause pain and fatigue, the commonest mechanism of TMD
    \item \textbf{Disc displacement:} the cushioning disc can slip forward or sideways, causing clicking, popping, and sometimes locking of the jaw
    \item \textbf{Inflammation and arthritis:} the joint lining can become inflamed, and osteoarthritis or inflammatory joint disease such as rheumatoid arthritis can affect the jaw joint as they do other joints
    \item \textbf{Injury:} a blow to the jaw, whiplash, prolonged mouth opening during dental treatment, or hard biting can strain the joint and muscles
    \item \textbf{Habitual overuse:} chewing gum constantly, biting nails, chewing ice, or resting the chin on the hand can aggravate the joint and muscles
    \item \textbf{Stress and anxiety:} psychological stress increases muscle tension, clenching, and pain perception, and is a major driver of flare-ups
    \item \textbf{Posture and neck problems:} poor head and neck posture is associated with jaw muscle tension, and neck and jaw problems frequently coexist
\end{itemize}

In many cases several factors combine: a person with an injury or with slight joint wear develops muscle overactivity, which produces pain, which increases stress, which increases clenching, and so a cycle is established. Breaking this cycle, rather than curing a single abnormality, is the goal of treatment.

\section*{Clinical Features}
\begin{itemize}
    \item \textbf{Jaw pain:} aching or throbbing pain in front of the ear, in the jaw, or in the cheek, worse with chewing, talking, or yawning
    \item \textbf{Pain on palpation:} tenderness of the joint in front of the ear and of the chewing muscles over the cheeks and temples
    \item \textbf{Clicking or popping:} audible or felt clicks, pops, or grating sounds when opening or closing the mouth; sounds without pain are usually harmless
    \item \textbf{Locking:} the jaw occasionally sticks open or closed, sometimes requiring the person to manipulate it to free it
    \item \textbf{Limited opening:} difficulty opening the mouth fully, with normal opening being roughly the width of three fingers
    \item \textbf{Difficulty chewing:} discomfort or reduced power with food, sometimes with a feeling that the bite is not fitting right
    \item \textbf{Headache:} morning headaches or tension-type headaches that spread from the temples, common when night-time grinding is a factor
    \item \textbf{Ear symptoms:} pain in or around the ear, a sense of fullness, or ringing without a clear ear cause
    \item \textbf{Neck pain:} associated neck and shoulder tension is common
    \item \textbf{Tooth wear:} a dentist may see flattened, worn teeth indicating grinding
\end{itemize}

\section*{Differential Diagnosis}
\begin{itemize}
    \item \textbf{Ear disease:} otitis media and eustachian tube problems cause ear pain and fullness that can be mistaken for TMD; a hearing check or examination of the ear settles it
    \item \textbf{Dental problems:} tooth decay, cracked teeth, impacted wisdom teeth, and dental abscesses cause jaw or cheek pain; a dental examination is essential
    \item \textbf{Sinusitis:} maxillary sinus infection causes cheek and tooth-area pain, often with nasal congestion
    \item \textbf{Trigeminal neuralgia:} brief, severe, electric-shock-like facial pain along a nerve branch, very different from the dull ache of TMD
    \item \textbf{Giant cell arteritis:} a serious inflammatory condition of scalp blood vessels in older people causing temple pain and, untreated, blindness; requires urgent assessment
    \item \textbf{Salivary gland disease:} stones or infection of the parotid gland cause swelling and pain around the jaw
    \item \textbf{Fibromyalgia and chronic pain syndromes:} widespread pain conditions can include jaw pain and should be considered when pain is generalised
    \item \textbf{Tumours (rare):} any persistent, unexplained swelling, numbness, or progressive limitation of mouth opening warrants imaging
\end{itemize}

\section*{When to Seek Medical Care}
Most TMD can be managed at home with simple measures, but professional assessment is worthwhile when symptoms persist.

Seek advice from a GP or dentist if you have:
\begin{itemize}
    \item Jaw pain that lasts more than a few weeks or is getting worse
    \item Locking of the jaw that does not free itself, or progressive difficulty opening the mouth
    \item Pain or noises associated with swelling, fever, or redness near the joint
    \item Difficulty or pain with eating that is causing weight loss or poor nutrition
    \item New headaches, ear symptoms, or facial pain without an obvious cause
    \item Clenching or grinding that is wearing down your teeth
\end{itemize}

\textbf{Contact your local emergency services or seek urgent care for:}
\begin{itemize}
    \item Inability to open the mouth after an injury, or sudden severe facial pain with swelling
    \item Facial pain with vision change, jaw claudication (pain on chewing), or scalp tenderness, which can indicate giant cell arteritis
    \item Swelling spreading rapidly, difficulty breathing or swallowing, or signs of severe infection
\end{itemize}

\section*{Diagnosis and Investigations}
The diagnosis of TMD is clinical, based on the history and a careful examination, and imaging is needed only when the diagnosis is uncertain or a complication is suspected.

\begin{itemize}
    \item \textbf{History:} questions about the pattern of pain, clicking, locking, chewing difficulty, grinding or clenching habits, stress, and past jaw injury
    \item \textbf{Physical examination:} measuring jaw opening, listening for joint sounds, and pressing gently over the joint and chewing muscles to locate tenderness
    \item \textbf{Dental examination:} checking the teeth for wear, cracked teeth, and bite problems, and ruling out dental causes of pain
    \item \textbf{Panoramic dental X-ray:} a single film of the whole mouth and jaws, useful to look at the joint position and exclude dental or bone problems
    \item \textbf{CBCT or MRI:} advanced imaging of the joint used when locking, suspected arthritis, or internal disc problems need confirmation; MRI is the best test for the disc itself
    \item \textbf{Blood tests:} inflammatory markers and, in the right setting, tests for rheumatoid arthritis or other systemic disease
    \item \textbf{Referral:} persistent or complex cases are referred to an oral medicine specialist, oral surgeon, or a TMD-focused physiotherapist
\end{itemize}

\section*{Management}
Conservative treatment is effective for the great majority, and the emphasis is on education, unloading the joint, and breaking the pain cycle.

\begin{itemize}
    \item \textbf{Education and reassurance:} understanding that most TMD settles and that surgery is rarely needed reduces anxiety, which itself reduces clenching
    \item \textbf{Soft diet and rest:} eat soft foods, cut food into small pieces, and avoid chewing gum, chewy foods, and wide biting during flare-ups
    \item \textbf{Heat and cold:} a warm pack over the jaw muscles, and a cold pack over an acutely painful joint, give relief
    \item \textbf{Jaw exercises:} gentle, guided opening and stretching exercises, often with physiotherapy, restore movement and reduce stiffness
    \item \textbf{Habit awareness:} keep the teeth apart and lips together at rest, avoid clenching during the day, and stop habits such as nail biting and pen chewing
    \item \textbf{Pain relief:} short-term paracetamol or an anti-inflammatory medicine such as ibuprofen, used as directed, helps settle acute pain
    \item \textbf{Sleep splint or bite splint:} a dentist-fitted acrylic splint worn at night reduces the damage and muscle loading from grinding, and is helpful in many people
    \item \textbf{Stress management:} relaxation, mindfulness, breathing exercises, and attention to sleep help break the stress--clenching cycle
    \item \textbf{Physiotherapy:} specific jaw and neck exercises, mobilisation, and postural advice from a trained physiotherapist are effective
    \item \textbf{Specialist treatments:} for severe cases, options include joint lavage, corticosteroid injection, and, rarely, surgery such as arthroscopy; these are reserved for people who do not respond to conservative care
\end{itemize}

\section*{Complications}
\begin{itemize}
    \item \textbf{Chronic pain:} untreated TMD can become a persistent facial pain problem, with the associated effects on sleep, eating, and mood
    \item \textbf{Locking:} the jaw can become stuck, occasionally requiring dental or hospital treatment to release it
    \item \textbf{Tooth damage:} continued grinding wears down enamel and can crack teeth or loosen dental work
    \item \textbf{Muscle strain and neck problems:} compensatory tension and posture changes can cause headaches and neck pain
    \item \textbf{Arthritis progression:} inflammatory or degenerative joint disease can worsen and limit jaw function over time
    \item \textbf{Psychological effects:} chronic jaw pain is associated with anxiety, poor sleep, and reduced quality of life if not addressed
\end{itemize}

\section*{Prognosis and Outcomes}
The outlook for TMD is generally very good. Most people improve substantially with conservative management within weeks to a few months, and even without treatment many episodes settle on their own once the provoking stress or behaviour passes. Clicking noises frequently persist even when pain resolves, and they rarely need treatment in themselves. The pattern is often one of flare-ups and quiet periods, so learning to recognise and respond early to a flare-up, with a soft diet, heat, and relaxation, is a practical and valuable skill. People who clench or grind at night may need a splint long term to protect the teeth. A small minority have persistent pain despite treatment; for them, a coordinated approach involving a dentist, a GP, physiotherapy, and often psychology gives the best chance of a good outcome. Surgical treatment is reserved for a small group with confirmed structural problems such as true locking or joint disease, and in that group outcomes are generally good.

\section*{Prevention and Self-Care}
\begin{itemize}
    \item Keep the jaw relaxed during the day: teeth apart, lips together, and avoid clenching when concentrating or stressed
    \item Eat a soft diet and chew evenly on both sides during flare-ups, and limit gum and chewy foods
    \item Practise relaxation and good sleep habits to reduce night-time grinding
    \item Use a warm pack on tight jaw muscles and a cold pack on a swollen, painful joint
    \item Avoid habits that stress the jaw, such as nail biting, pen chewing, ice chewing, and resting the chin on the hand
    \item Improve head and neck posture, particularly when working at a desk or using a phone
    \item See a dentist if you grind your teeth or if tooth wear is evident; a night splint protects the teeth and the joint
    \item Seek help early if pain persists beyond a few weeks, rather than waiting for the joint to worsen
\end{itemize}

\section*{Sources and Further Reading}
The following organisations provide reliable information about temporomandibular joint dysfunction and jaw pain.

\begin{itemize}
    \item Australian Dental Association. \textit{Temporomandibular joint (TMJ) dysfunction.} https://www.ada.org.au/
    \item healthdirect Australia. \textit{TMJ dysfunction and jaw pain.} https://www.healthdirect.gov.au/
    \item NHS. \textit{TMJ disorders.} https://www.nhs.uk/
    \item Mayo Clinic. \textit{TMJ disorders.} https://www.mayoclinic.org/
    \item Better Health Channel (Victoria). \textit{Jaw problems and TMJ dysfunction.} https://www.betterhealth.vic.gov.au/
    \item Australian Department of Health and Aged Care. \textit{Dental and oral health resources.} https://www.health.gov.au/
\end{itemize}

\clearpage
